\documentclass[11pt]{article}
\usepackage[T1]{fontenc}
\usepackage{amsmath,amssymb,amsthm}
\usepackage[margin=1in]{geometry}
\usepackage[hidelinks]{hyperref}
\newtheorem{theorem}{Theorem}
\newtheorem{lemma}{Lemma}
\newtheorem{proposition}{Proposition}
\newcommand{\R}{\mathbb R}
\newcommand{\C}{\mathbb C}
\newcommand{\ip}[2]{\langle #1,#2\rangle}
\newcommand{\norm}[1]{\lVert #1\rVert}
\title{Nonunique nearest stable polynomials near repeated roots}
\author{Model-assisted research contribution}
\date{28 October 2026}
\begin{document}
\maketitle

\begin{abstract}
For the Euclidean norm of the nonleading coefficients, the closed set of real monic stable quartics is not prox-regular at \(z^4\). Consequently, arbitrarily near \(z^4\) there are real monic quartics with at least two nearest real monic stable quartics. We separately prove the analogous conclusion for complex monic quadratics near \(z^2\). The proofs use explicit stable boundary paths and proximal normals. These are negative results for the monic coefficient-space version of the nearest-stable-polynomial question; no conclusion about finite-subgradient prox-regularity of the abscissa function is asserted.
\end{abstract}

\section{Question, conventions, and scope}

The nearest-stable-polynomial question in Lewis's 2005 report \cite[p.~106, Question~1]{owr} asks about uniqueness near a monomial. The preceding discussion uses real polynomials, while the final question does not restate a coefficient field or norm. Lewis's earlier monic formulation \cite[Section~8]{lewis} uses complex coefficients and the real Euclidean structure making the coefficient basis and its imaginary multiples orthonormal. We treat both real and complex monic coefficient spaces explicitly.

For \(\mathbb F=\R\) or \(\C\), let
\[
 S_{n,\mathbb F}
 =\left\{c=(c_1,\ldots,c_n)\in\mathbb F^n:
       z^n+c_1z^{n-1}+\cdots+c_n
       \text{ has all roots in }\{\operatorname{Re}z\leq0\}\right\}.
\]
In complex space the inner product is
\(\ip{c}{d}=\operatorname{Re}\sum_j\overline{c_j}d_j\).
The norm is its Euclidean norm. A polynomial is stable here in the closed-half-plane sense, and zero coefficients represent \(z^n\). Write \(P_S(y)\) for the set of global nearest points to \(y\) in \(S\).

\begin{theorem}\label{thm:main}
For every \(\delta>0\):
\begin{enumerate}
\item There is \(y\in\R^4\), \(\norm{y}<\delta\), such that
      \(P_{S_{4,\R}}(y)\) contains at least two distinct points.
\item There is \(y\in\C^2\), \(\norm{y}<\delta\), such that
      \(P_{S_{2,\C}}(y)\) contains at least two distinct points.
\end{enumerate}
In each case all nearest points have norm at most \(2\norm{y}\).
The sets \(S_{4,\R}\) and \(S_{2,\C}\) are not prox-regular at zero.
\end{theorem}

The quartic stability boundary and its singularities are classical; for example, Kirillov and Overton \cite[Section~1]{kirillov} discuss the Routh--Hurwitz/Bottema geometry in a constant-term-normalized slice. Neither that boundary formula nor root sensitivity is claimed as a new result here. The contribution is the explicit normal-support obstruction at the monomial and its exact nearest-point consequence. The archived attempt for the present problem only treats the real monic quadratic case, whose stable coefficient set is a quadrant. The real quartic conclusion addresses the same projection question beyond that convex case.

The cited historical sources and targeted primary-literature searches did not yield an equivalent negative projection theorem in the material checked. This is a bounded comparison, not a comprehensive priority certification or a claim that the question's current status has been exhaustively surveyed.

\section{Normal vectors and projection uniqueness}

We recall the precise finite-dimensional facts needed from Poliquin, Rockafellar and Thibault \cite{prt}. A vector \(v\) is proximal normal to a closed set \(S\) at \(p\in S\) if \(p\in P_S(p+\tau v)\) for some \(\tau>0\). Equivalently, for some \(K\geq0\),
\begin{equation}\label{eq:quadratic}
 \ip{v}{q-p}\leq K\norm{q-p}^2
 \qquad(q\in S).
\end{equation}
Indeed, expanding the squared distance proves necessity and proves sufficiency for a sufficiently small positive \(\tau\). The limiting normal cone \(N_S(p)\) consists of limits of proximal normals at points converging to \(p\), together with zero.

The definition of prox-regularity at \(\bar p\) for
\(\bar v\in N_S(\bar p)\), in the form of \cite[Definition~1.1]{prt}, requires constants \(\eta,\rho>0\) such that, for nearby \(p\in S\) and \(v\in N_S(p)\) near \(\bar v\), \(p\) is the unique nearest point in \(S\cap B(\bar p,\eta)\) to \(p+\rho^{-1}v\). In particular it requires
\begin{equation}\label{eq:uniform}
 \ip{v}{q-p}\leq \frac{\rho}{2}\norm{q-p}^2
 \qquad(q\in S\cap B(\bar p,\eta)).
\end{equation}
Prox-regularity at a point requires this property for every limiting normal there.

We invoke the following exact consequence of
\cite[Theorem~1.3(a),(k)]{prt}: for a closed subset of a finite-dimensional Euclidean space, prox-regularity at \(\bar p\) is equivalent to single-valuedness of the global metric projection on some neighborhood of \(\bar p\). Its additional relative weak-closure hypothesis is automatic in finite dimension. We will disprove prox-regularity using \eqref{eq:uniform}.

\begin{lemma}\label{lem:smooth}
Suppose a closed set \(S\) agrees near \(p=(x_0,f(x_0))\) with the hypograph of a \(C^2\) function \(f\). Then \((-\nabla f(x_0),1)\) is a proximal normal at \(p\). For an epigraph the corresponding normal is \((\nabla f(x_0),-1)\).
\end{lemma}
\begin{proof}
For the hypograph, write \(q=(x,u)\). Taylor's inequality and \(u\leq f(x)\) imply
\[
 \ip{(-\nabla f(x_0),1)}{q-p}
 \leq f(x)-f(x_0)-\nabla f(x_0)\cdot(x-x_0)
 \leq K\norm{q-p}^2
\]
for \(q\in S\) in a sufficiently small ball about \(p\).
For other \(q\in S\), whose distance from \(p\) is at least some \(r>0\), Cauchy--Schwarz gives the same inequality with
\(K\) enlarged to at least
\(\norm{(-\nabla f(x_0),1)}/r\).
Thus \eqref{eq:quadratic} holds globally. The epigraph proof reverses the signs.
\end{proof}

\section{Real quartics}

Write a real monic quartic as
\[
 Q(z)=z^4+az^3+bz^2+cz+d.
\]
For \(0<\varepsilon\leq1/2\), put
\[
 p_\varepsilon=
 \left(\varepsilon^2,\varepsilon,
       \frac{\varepsilon^3}{4},\frac{3\varepsilon^2}{16}\right).
\]
It is stable because
\begin{equation}\label{eq:factor}
 Q_\varepsilon(z)=
 (z^2+\varepsilon/4)
 (z^2+\varepsilon^2z+3\varepsilon/4).
\end{equation}
The first factor has a simple imaginary conjugate pair; both roots of the second have strictly negative real part and are distinct, since its discriminant is \(\varepsilon^4-3\varepsilon<0\). The two factors share no root, since evaluating the second factor at either root of the first gives
\(\varepsilon/2\mathbin{\pm}i\varepsilon^2\sqrt{\varepsilon/4}\ne0\).

\begin{lemma}\label{lem:quarticboundary}
Near \(p_\varepsilon\), the stable real quartics form exactly the hypograph
\[
 d\leq f(a,b,c),\qquad
 f(a,b,c)=\frac{bc}{a}-\frac{c^2}{a^2}.
\]
\end{lemma}
\begin{proof}
The simple roots vary smoothly near this point by the implicit function theorem; the strictly left-half-plane pair stays there. The other two roots remain a conjugate pair near the imaginary pair. If one is \(i\sqrt y\), with \(y>0\), the imaginary and real parts of \(Q(i\sqrt y)=0\) give
\[
 y=c/a,\qquad d=by-y^2=f(a,b,c).
\]
Conversely, on this graph,
\[
 Q(z)=(z^2+c/a)(z^2+az+b-c/a).
\]
In a neighborhood of \(p_\varepsilon\), \(a>0\), \(c/a>0\), and \(b-c/a>0\); the second quadratic is strictly stable.
To determine the side, at an imaginary root \(z=i\sqrt y\) one has
\[
 Q'(i\sqrt y)=-2ay+2i\sqrt y\,(b-2y),\qquad
 \operatorname{Re}\left(-\frac{1}{Q'(i\sqrt y)}\right)>0.
\]
This is the derivative of the active root's real part with respect to \(d\). It is nonzero, so the graph separates the stable and unstable sides locally, and increasing \(d\) gives the unstable side. The claimed hypograph follows.
\end{proof}

At \(p_\varepsilon\),
\[
 f_a=-\frac18,\qquad f_b=\frac{\varepsilon}{4},
 \qquad f_c=\frac{1}{2\varepsilon}.
\]
By Lemma~\ref{lem:smooth}, the vector
\[
 n_\varepsilon=
 \left(\frac18,-\frac{\varepsilon}{4},
              -\frac{1}{2\varepsilon},1\right)
\]
is proximal normal. Its normalization
\(v_\varepsilon=n_\varepsilon/\norm{n_\varepsilon}\)
converges to \(v_0=(0,0,-1,0)\), so \(v_0\in N_{S_{4,\R}}(0)\).

Keep \(b=\varepsilon\), \(c=\varepsilon^3/4\) fixed, set \(a=\varepsilon^2\), and choose \(h=\varepsilon^4\). Define
\[
 q_\varepsilon=(a+h,b,c,f(a+h,b,c)).
\]
This point is stable by the factorization in Lemma~\ref{lem:quarticboundary}; its positive quadratic constant is at least \(3\varepsilon/4\). It converges to zero. Direct subtraction gives the exact positive support gap
\begin{align}
 \ip{n_\varepsilon}{q_\varepsilon-p_\varepsilon}
 &=f(a+h,b,c)-f(a,b,c)-f_a(a,b,c)h \notag\\
 &=\frac{h^2(\varepsilon^2+2h)}
          {16(\varepsilon^2+h)^2}
 \geq\frac{h^2}{64\varepsilon^2}.             \label{eq:gap}
\end{align}
Here \(h\leq\varepsilon^2\). On \(A\in[a,a+h]\), one has
\[
 -\frac14\leq
 \frac{\partial f}{\partial A}
 =-\frac{bc}{A^2}+\frac{2c^2}{A^3}\leq0.
\]
The mean value theorem therefore implies
\(\norm{q_\varepsilon-p_\varepsilon}^2\leq2h^2\).
Also, for \(0<\varepsilon\leq1/2\),
\[
 \norm{n_\varepsilon}^2
 =\frac{65}{64}+\frac{\varepsilon^2}{16}
                    +\frac{1}{4\varepsilon^2}
 \leq\frac{1}{\varepsilon^2}.
\]
Combining these estimates with \eqref{eq:gap} yields
\begin{equation}\label{eq:divergence}
 \frac{\ip{v_\varepsilon}{q_\varepsilon-p_\varepsilon}}
      {\norm{q_\varepsilon-p_\varepsilon}^2}
 \geq\frac{1}{128\varepsilon}\longrightarrow+\infty.
\end{equation}
Both points approach zero and \(v_\varepsilon\to v_0\). Thus no fixed \(\rho,\eta\) can satisfy \eqref{eq:uniform} at \((0,v_0)\). This disproves prox-regularity of \(S_{4,\R}\) at zero.

\section{Complex quadratics}

Now work independently in \(\C^2\), with
\(Q(z)=z^2+az+b\). Take
\[
 p_\varepsilon=(\varepsilon^2,\varepsilon^2+i\varepsilon^3),
 \qquad
 Q_\varepsilon(z)=(z+i\varepsilon)(z+\varepsilon^2-i\varepsilon).
\]
The roots are distinct, one on the imaginary axis and one strictly in the left half-plane. Use real orthonormal coordinates
\[
 (x,s,t,u)=(\operatorname{Re}a,\operatorname{Im}a,
            \operatorname{Im}b,\operatorname{Re}b).
\]
In a neighborhood of \(p_\varepsilon\), \(x>0\). An imaginary active root \(iy\) satisfies
\[
 y=-t/x,\qquad u=y^2+sy
       =f(x,s,t):=\frac{t^2}{x^2}-\frac{st}{x}.
\]
Conversely these equations give a root \(iy\), with the other root having real part \(-x<0\).
At \(p_\varepsilon\), the derivative of the active root's real part with respect to \(u\) is
\[
 \operatorname{Re}\left(-\frac{1}{\varepsilon^2-2i\varepsilon}\right)
 =-\frac{1}{\varepsilon^2+4}<0.
\]
The same simple-root implicit function argument as above shows that the local stable set is exactly the epigraph \(u\geq f(x,s,t)\).

Its outward proximal normal, by Lemma~\ref{lem:smooth}, is
\[
 w_\varepsilon=(\nabla f,-1)
       =(-2,-\varepsilon,2/\varepsilon,-1).
\]
Since \(p_\varepsilon\to0\) and
\(w_\varepsilon/\norm{w_\varepsilon}\to v=(0,0,1,0)\),
the vector \(v\) belongs to \(N_{S_{2,\C}}(0)\).
But
\[
 \ip{v}{p_\varepsilon}=\varepsilon^3,\qquad
 \norm{p_\varepsilon}^2=2\varepsilon^4+\varepsilon^6.
\]
For every fixed \(\rho>0\),
\[
 \norm{p_\varepsilon-\rho^{-1}v}^2-\norm{\rho^{-1}v}^2
 =2\varepsilon^4+\varepsilon^6-2\rho^{-1}\varepsilon^3<0
\]
for sufficiently small \(\varepsilon>0\). Thus zero is not even a local nearest point to \(\rho^{-1}v\) in the stable set. Definition~1.1 of \cite{prt}, applied with the point zero and normal \(v\) themselves, fails for every \(\rho\). Hence \(S_{2,\C}\) is not prox-regular at zero.

\section{Nearest-point conclusion and limitations}

For each fixed \(n\), \(S_{n,\mathbb F}\) is closed. Indeed, bounded monic coefficients give a uniform bound on all roots. From any convergent coefficient sequence of stable polynomials one can select a convergent subsequence of their ordered root tuples; the limiting factorization has only roots in the closed left half-plane. Also \(0\in S_{n,\mathbb F}\). In finite dimension a bounded minimizing sequence therefore has a convergent subsequence attaining the distance, so every global projection is nonempty.

The two failures of prox-regularity, together with
\cite[Theorem~1.3(a),(k)]{prt}, show that neither projection is single-valued on any neighborhood of zero. Nonemptiness turns this into the two asserted nonuniqueness statements in Theorem~\ref{thm:main}. If \(q\in P_S(y)\), then
\[
 \norm{q-y}\leq\norm{y},\qquad \norm{q}\leq2\norm{y},
\]
because zero is feasible. Thus the competing nearest stable polynomials also approach the monomial.

The complex quadratic conclusion is proved in the full four-real-dimensional coefficient space, not inferred from a real subset. The quartic conclusion concerns real coefficients and their own Euclidean projection. No implication between the two ambient spaces is used. We do not classify all degrees, change the norm, vary the leading coefficient, or assert an explicit closed formula for the centers with nonunique projection. Their existence follows from the precise finite-dimensional projection theorem. Strict stability would give an open feasible set and is a different nearest-point question.

Finally, these are statements about the stable set, or equivalently its indicator function. They do not prove failure of prox-regularity of the polynomial abscissa function at a finite subgradient, and they do not conflate projection onto a zero sublevel set with projection onto that function's epigraph.

\section*{Provenance, assistance, and licence}

The problem is
\texttt{prb\_4c50f9b914a642b2bc37e90aff6d8120}, statement revision
\texttt{stm\_547a66e8ccd74186be2bf7d41da2b035}, checked with its full archived attempt at We-Solve/research commit
\texttt{c4ef184d1d392f70fed32c5e862d7f24d11e8401}.
The immutable statement is
\url{https://github.com/We-Solve/research/blob/c4ef184d1d392f70fed32c5e862d7f24d11e8401/problems/4c/prb_4c50f9b914a642b2bc37e90aff6d8120/statement.tex}.

The derivations and manuscript were produced by OpenAI
\texttt{gpt-6.1-sol}, reasoning setting \texttt{ultra}, as explicitly configured for research agent e; exact serving revision and time budget were not exposed. Research agent f, under the same model/settings, independently challenged the proofs and suggested the bound keeping all nearest points near the monomial. Agent a assisted only with a bounded primary-literature overlap search. The coordinating Codex agent, based on GPT-6 with exact revision/settings unknown, supplied scope challenges and literature pointers. No human mathematical authorship or independent human proof verification is claimed.

The argument is analytic; no numerical experiment or formal proof checker is claimed. Independent mathematical and source-format checks are recorded separately. No local TeX compiler was available, so compilation and PDF rendering were not performed; source-format validation is not compilation. Original exposition and arguments are offered under CC BY 4.0; cited source material retains its own attribution and licence.

\begin{thebibliography}{9}
\bibitem{owr}
A.~S. Lewis,
\emph{Local Structure and Algorithms in Nonsmooth Optimization},
Oberwolfach Report 2/2005, pp.~104--106, especially Question~1.
\url{https://ems.press/content/serial-article-files/45976}.

\bibitem{lewis}
A.~S. Lewis,
\emph{The mathematics of eigenvalue optimization},
Mathematical Programming, Series B \textbf{97} (2003), 155--176,
Section~8, printed pp.~166--168.
\url{https://doi.org/10.1007/s10107-003-0441-3}.
Author manuscript:
\url{https://optimization-online.org/wp-content/uploads/2003/04/640.pdf}.

\bibitem{prt}
R.~A. Poliquin, R.~T. Rockafellar, and L.~Thibault,
\emph{Local differentiability of distance functions},
Transactions of the American Mathematical Society
\textbf{352} (2000), 5231--5249,
Definition~1.1, Proposition~1.2, and Theorem~1.3.
\url{https://doi.org/10.1090/S0002-9947-00-02550-2}.
Author copy:
\url{https://sites.math.washington.edu/~rtr/papers/rtr170-LocalDiffDistance.pdf}.

\bibitem{kirillov}
O.~N. Kirillov and M.~L. Overton,
\emph{Robust stability at the Swallowtail singularity},
Frontiers in Physics \textbf{1} (2013), Article~24,
Section~1.
\url{https://doi.org/10.3389/fphy.2013.00024}.
\end{thebibliography}
\end{document}
